A continuous linear map $u: L \to M$ between topological vector spaces over $\mathbb{K}$ is called a \textit{topological homomorphism} (or simply a homomorphism) if for each open subset $G \subset L$, the image $u(G)$ is an open subset of $u(L)$ with respect to the topology induced by $M$.

Every linear map $u: L \to M$ admits a canonical decomposition

$$L \xrightarrow{\phi} L/N \xrightarrow{u_0} u(L) \xrightarrow{\psi} M$$

where $N = u^{-1}(0)$ is the null space, $\phi: L \to L/N$ is the canonical quotient map, $\psi: u(L) \to M$ is the canonical imbedding, and $u_0$ is the bijective map associated with $u$ sending each coset $\hat{x} \in L/N$ to $u(x)$. Thus $u = \psi \circ u_0 \circ \phi$.
